\section{Flujo eléctrico}

En el estudio del campo eléctrico hemos aprendido a determinar la magnitud y la
dirección de \(\vec E\) en diferentes puntos del espacio.

Ahora introduciremos una magnitud que permite describir cuánto campo eléctrico
atraviesa una superficie.

Esta magnitud se denomina \textbf{flujo eléctrico}.

La idea puede entenderse inicialmente mediante las líneas de campo. Si muchas líneas
atraviesan una superficie, asociamos a esa situación un flujo eléctrico grande. Si pocas
líneas la atraviesan, el flujo es menor.

Sin embargo, las líneas de campo son solamente una representación gráfica. El flujo
eléctrico se define matemáticamente a partir del campo eléctrico y de la geometría de
la superficie.


\subsection{Campo uniforme perpendicular a una superficie}

Consideremos una superficie plana de área \(A\) colocada en un campo eléctrico
uniforme de magnitud \(E\).

Si el campo atraviesa perpendicularmente la superficie, el flujo eléctrico es

\[
\boxed{\Phi_E=EA}.
\]

El flujo aumenta si aumenta la intensidad del campo o si aumenta el área atravesada.

\begin{idea}

El flujo eléctrico no mide simplemente el valor del campo eléctrico.

También depende del tamaño y de la orientación de la superficie que el campo
atraviesa.

\end{idea}


\subsection{Vector área}

Para describir la orientación de una superficie plana asociamos a ella un
\textbf{vector área} \(\vec A\).

Su magnitud es igual al área de la superficie:

\[
|\vec A|=A,
\]

y su dirección es perpendicular a la superficie.

Una superficie plana posee dos posibles direcciones normales. En una superficie
abierta puede elegirse cualquiera de ellas siempre que se mantenga la elección de
forma consistente.

En una superficie cerrada existe una convención específica: el vector área apunta
siempre hacia el exterior.


\subsection{Superficie inclinada respecto del campo}

Supongamos ahora que el campo eléctrico forma un ángulo \(\theta\) con la dirección
normal a la superficie.

El flujo eléctrico es

\begin{equation}
\boxed{\Phi_E=EA\cos\theta}.
\label{eq:flujo_superficie_plana}
\end{equation}

El ángulo \(\theta\) se mide entre el campo eléctrico y la \textbf{normal a la superficie},
no entre el campo y el plano de la superficie.

\begin{center}
\begin{tikzpicture}[>=Latex,scale=1.0]

\draw[thick] (-2,-0.6) -- (1.8,0.5) -- (3.2,1.8) -- (-0.6,0.7) -- cycle;

\draw[->,very thick] (-3,0.6) -- (3.7,0.6) node[right] {$\vec E$};

\coordinate (C) at (0.6,0.55);
\draw[->,very thick] (C) -- (1.8,2.9) node[above] {normal};

\draw (1.15,0.55) arc[start angle=0,end angle=63,radius=0.55];
\node at (1.55,1.05) {$\theta$};

\node at (0,-0.55) {superficie de área \(A\)};

\end{tikzpicture}
\end{center}

La cantidad

\[
A\cos\theta
\]

puede interpretarse como el área efectiva perpendicular al campo.


\subsection{Casos particulares}

Si

\[
\theta=0^\circ,
\]

el campo es perpendicular a la superficie y paralelo a su normal. Entonces,

\[
\Phi_E=EA.
\]

El flujo tiene su valor máximo positivo.

Si

\[
\theta=90^\circ,
\]

el campo es paralelo a la superficie. Entonces,

\[
\boxed{\Phi_E=0}.
\]

En este caso el campo no atraviesa la superficie.

Si

\[
\theta=180^\circ,
\]

el campo es opuesto a la normal:

\[
\boxed{\Phi_E=-EA}.
\]

El signo negativo indica que el campo atraviesa la superficie en sentido contrario al
vector área elegido.


\begin{advertencia}

Un flujo igual a cero no implica necesariamente que

\[
E=0.
\]

Puede existir un campo eléctrico distinto de cero y, sin embargo, ser paralelo a la
superficie, de manera que no la atraviese.

\end{advertencia}


\subsection{Unidad del flujo eléctrico}

Como

\[
\Phi_E=EA,
\]

la unidad SI es

\[
[\Phi_E]
=
\left(\si{\newton\per\coulomb}\right)
\left(\si{\metre\squared}\right).
\]

Por tanto,

\[
\boxed{
[\Phi_E]=\si{\newton\metre\squared\per\coulomb}
}.
\]


\subsection{Ejemplo: flujo a través de una superficie plana}

\begin{ejemplo}

Una superficie plana de área

\[
A=\SI{0.20}{\metre\squared}
\]

se encuentra en un campo eléctrico uniforme de magnitud

\[
E=\SI{500}{\newton\per\coulomb}.
\]

La normal a la superficie forma un ángulo de \(60^\circ\) con el campo.

Determine el flujo eléctrico.

\end{ejemplo}

Aplicamos

\[
\Phi_E=EA\cos\theta.
\]

Sustituyendo,

\[
\Phi_E=(500)(0.20)\cos60^\circ.
\]

Como

\[
\cos60^\circ=0.50,
\]

obtenemos

\[
\boxed{
\Phi_E=\SI{50}{\newton\metre\squared\per\coulomb}
}.
\]

Si en lugar de \(60^\circ\) con la normal nos hubieran dado el ángulo con el plano de
la superficie, primero habría sido necesario obtener el ángulo complementario.


\section{Flujo a través de una superficie no uniforme}

La expresión

\[
\Phi_E=EA\cos\theta
\]

sólo puede utilizarse directamente cuando el campo es uniforme sobre toda la superficie
y la orientación de ésta es constante.

Si la superficie es curva o el campo cambia de un punto a otro, debemos dividirla
conceptualmente en pequeños elementos de área \(dA\).

Cada elemento puede considerarse aproximadamente plano.

El flujo infinitesimal es

\[
d\Phi_E=E\cos\theta\,dA.
\]

Por tanto, el flujo total se obtiene integrando:

\begin{equation}
\boxed{
\Phi_E=\int_S \vec E\cdot d\vec A
}.
\label{eq:flujo_superficie_general}
\end{equation}

De forma escalar,

\[
\boxed{
\Phi_E=\int_S E\cos\theta\,dA
}.
\]

Aquí \(\theta\) puede cambiar de un punto a otro.


\begin{idea}

La integral de flujo suma únicamente la componente del campo perpendicular a cada
elemento de superficie.

La componente paralela a la superficie no contribuye al flujo.

\end{idea}


\section{Superficies abiertas y cerradas}

Una superficie puede ser abierta o cerrada.

Una superficie plana, un disco o una hoja son ejemplos de superficies abiertas.

Una esfera, un cubo o un cilindro con sus dos tapas son ejemplos de superficies cerradas.

En una superficie cerrada se utiliza el símbolo

\[
\oint
\]

para indicar que la integración se realiza sobre toda la superficie.

El flujo neto es

\begin{equation}
\boxed{
\Phi_E=\oint_S\vec E\cdot d\vec A
}.
\label{eq:flujo_superficie_cerrada}
\end{equation}

Para una superficie cerrada, \(d\vec A\) apunta siempre hacia afuera.


\subsection{Flujo positivo y negativo}

En una superficie cerrada:

\[
\boxed{\text{campo saliente}\Longrightarrow d\Phi_E>0}
\]

y

\[
\boxed{\text{campo entrante}\Longrightarrow d\Phi_E<0}.
\]

Por ello, el flujo neto es la diferencia entre las contribuciones salientes y entrantes.


\subsection{Flujo neto igual a cero}

Puede ocurrir que parte del campo entre a una superficie cerrada y otra parte salga.

Si las contribuciones positivas y negativas se cancelan,

\[
\Phi_E=0.
\]

Esto no significa que el campo sea cero sobre la superficie.

\begin{errorfrecuente}

No confundir

\[
\boxed{\Phi_E=0}
\]

con

\[
\boxed{E=0}.
\]

Son afirmaciones diferentes.

Un flujo neto nulo significa únicamente que las contribuciones entrantes y salientes
se compensan.

\end{errorfrecuente}


\section{Carga encerrada}

Para aplicar la ley de Gauss es necesario distinguir entre la carga
\textbf{encerrada} y las cargas que se encuentran fuera de una superficie cerrada.

Denotaremos por

\[
Q_{\text{enc}}
\]

la carga neta total encerrada.

Si dentro existen varias cargas,

\[
\boxed{Q_{\text{enc}}=\sum_iq_i}.
\]

Debe realizarse una suma algebraica, teniendo en cuenta sus signos.

Por ejemplo, si una superficie encierra

\[
+\SI{5}{\micro\coulomb},
\qquad
-\SI{2}{\micro\coulomb},
\qquad
+\SI{1}{\micro\coulomb},
\]

entonces

\[
Q_{\text{enc}}
=
5-2+1
=
+\SI{4}{\micro\coulomb}.
\]


\section{Ley de Gauss}

Existe una relación fundamental entre el flujo eléctrico neto a través de una superficie
cerrada y la carga neta que dicha superficie encierra.

\begin{principio}[Ley de Gauss]

El flujo eléctrico neto a través de cualquier superficie cerrada es igual a la carga neta
encerrada dividida entre la permitividad eléctrica del vacío.

\end{principio}

Matemáticamente,

\begin{equation}
\boxed{
\oint_S\vec E\cdot d\vec A
=
\frac{Q_{\text{enc}}}{\varepsilon_0}
}.
\label{eq:ley_gauss}
\end{equation}

Como el miembro izquierdo es el flujo eléctrico neto,

\begin{equation}
\boxed{
\Phi_E=\frac{Q_{\text{enc}}}{\varepsilon_0}
}.
\label{eq:gauss_flujo}
\end{equation}

La permitividad eléctrica del vacío es aproximadamente

\[
\boxed{
\varepsilon_0\approx
\SI{8.854e-12}{\coulomb\squared\per\newton\metre\squared}
}.
\]


\subsection{Significado físico}

La ley de Gauss establece que el flujo eléctrico neto a través de una superficie cerrada
depende exclusivamente de la carga neta encerrada.

Por tanto,

\[
Q_{\text{enc}}>0
\Longrightarrow
\Phi_E>0,
\]

\[
Q_{\text{enc}}<0
\Longrightarrow
\Phi_E<0,
\]

y

\[
Q_{\text{enc}}=0
\Longrightarrow
\Phi_E=0.
\]

Una carga positiva encerrada produce un exceso neto de flujo saliente.

Una carga negativa encerrada produce un exceso neto de flujo entrante.


\section{Las cargas externas y el flujo neto}

Una carga situada fuera de una superficie cerrada puede producir un campo eléctrico
sobre dicha superficie.

Sin embargo, su contribución al flujo \textbf{neto} es cero.

Las líneas asociadas a una carga externa que entran en la superficie deben volver a salir.

La entrada produce flujo negativo y la salida flujo positivo, de modo que el balance es

\[
\boxed{\Phi_{E,\text{carga externa}}=0}.
\]

\begin{advertencia}

Una carga externa puede modificar considerablemente el campo eléctrico sobre una
superficie gaussiana.

Lo que no modifica es el flujo neto asociado a la carga encerrada.

La ley de Gauss relaciona

\[
\oint_S\vec E\cdot d\vec A
\]

con \(Q_{\text{enc}}\), no con la carga total existente en todo el espacio.

\end{advertencia}


\section{La forma y el tamaño de la superficie cerrada}

Según la ley de Gauss,

\[
\Phi_E=\frac{Q_{\text{enc}}}{\varepsilon_0}.
\]

Por tanto, mientras la misma carga permanezca encerrada, el flujo total no depende:

\begin{itemize}
    \item del tamaño de la superficie;
    \item de su forma;
    \item de la posición de la carga dentro de la superficie.
\end{itemize}

La distribución local del flujo sí puede cambiar, pero el flujo neto permanece igual.


\subsection{Ejemplo conceptual}

\begin{ejemplo}

Una carga puntual \(+q\) está dentro de una superficie cerrada.

¿Qué sucede con el flujo total si:

\begin{enumerate}
    \item se duplica el tamaño de la superficie;
    \item la esfera se reemplaza por un cubo;
    \item la carga se desplaza dentro de la superficie sin salir de ella?
\end{enumerate}

\end{ejemplo}

En los tres casos la carga encerrada continúa siendo \(+q\).

Por tanto,

\[
\boxed{\Phi_E=\frac{q}{\varepsilon_0}}
\]

permanece constante.

El campo puede cambiar de un punto a otro de la nueva superficie, pero el flujo neto
no cambia.


\section{Superficie gaussiana}

Una \textbf{superficie gaussiana} es una superficie cerrada imaginaria elegida para
aplicar la ley de Gauss.

No necesita coincidir con una superficie física.

\begin{idea}

La superficie gaussiana es una herramienta matemática.

No es una barrera, una membrana ni una superficie material.

\end{idea}

La ley de Gauss puede aplicarse a cualquier superficie cerrada.

Sin embargo, si queremos utilizarla para calcular el campo eléctrico, no cualquier
superficie resulta conveniente.


\section{Ley de Gauss y simetría}

La ley de Gauss es válida para cualquier distribución de carga.

Pero para despejar fácilmente \(E\) de

\[
\oint_S\vec E\cdot d\vec A
=
\frac{Q_{\text{enc}}}{\varepsilon_0},
\]

debemos conocer suficiente información sobre el campo en la superficie elegida.

Esto ocurre especialmente cuando existe:

\begin{itemize}
    \item simetría esférica;
    \item simetría cilíndrica;
    \item simetría plana.
\end{itemize}

Una buena superficie gaussiana debe aprovechar esa simetría.


\subsection{Qué buscamos al elegir una superficie gaussiana}

Conviene escoger una superficie donde ocurra una o varias de las siguientes situaciones:

\begin{enumerate}
    \item la magnitud \(E\) es constante sobre una parte de la superficie;
    \item el campo es perpendicular a la superficie, de modo que
    \[
    \cos\theta=1;
    \]
    \item el campo es paralelo a la superficie, de modo que
    \[
    \cos\theta=0;
    \]
    \item alguna región posee campo nulo.
\end{enumerate}

Estas condiciones permiten simplificar la integral de flujo.


\begin{advertencia}

La superficie gaussiana no se elige porque ``se parece'' simplemente al dibujo del
problema.

Debe elegirse de forma que refleje la \textbf{simetría de la distribución de carga y
del campo eléctrico}.

\end{advertencia}


\section{Relación entre la ley de Coulomb y la ley de Gauss}

Consideremos una carga puntual positiva \(Q\) y una superficie gaussiana esférica
de radio \(r\) centrada en ella.

Por simetría esférica:

\begin{itemize}
    \item el campo es radial;
    \item tiene la misma magnitud en todos los puntos de la esfera;
    \item es perpendicular a la superficie.
\end{itemize}

Por tanto,

\[
\oint_S\vec E\cdot d\vec A
=
E\oint_SdA.
\]

El área de la esfera es

\[
A=4\pi r^2.
\]

Entonces,

\[
\Phi_E=E(4\pi r^2).
\]

La ley de Gauss establece

\[
E(4\pi r^2)=\frac{Q}{\varepsilon_0}.
\]

Despejando,

\[
E=\frac{1}{4\pi\varepsilon_0}\frac{Q}{r^2}.
\]

Como

\[
k=\frac{1}{4\pi\varepsilon_0},
\]

obtenemos

\[
\boxed{E=k\frac{Q}{r^2}}.
\]

Para la magnitud,

\[
\boxed{E=k\frac{|Q|}{r^2}}.
\]

Por tanto, para una carga puntual, la ley de Gauss conduce al mismo resultado obtenido
a partir de la ley de Coulomb.


\begin{interpretacion}

El flujo a través de la esfera es

\[
\Phi_E=E(4\pi r^2).
\]

El campo disminuye como

\[
\frac{1}{r^2},
\]

pero el área de la esfera aumenta como

\[
r^2.
\]

Ambos efectos se compensan, por lo que el flujo total es independiente del radio:

\[
\boxed{\Phi_E=\frac{Q}{\varepsilon_0}}.
\]

\end{interpretacion}


\section{Procedimiento para aplicar la ley de Gauss}

Cuando se quiere determinar un campo eléctrico mediante la ley de Gauss, conviene
seguir un procedimiento sistemático.

\begin{enumerate}
    \item Identificar la simetría de la distribución de carga.
    \item Determinar cualitativamente la dirección del campo.
    \item Elegir una superficie gaussiana compatible con la simetría.
    \item Analizar qué partes de la superficie aportan flujo.
    \item Calcular la carga encerrada \(Q_{\text{enc}}\).
    \item Aplicar
    \[
    \oint_S\vec E\cdot d\vec A
    =
    \frac{Q_{\text{enc}}}{\varepsilon_0}.
    \]
    \item Simplificar la integral utilizando la simetría.
    \item Despejar \(E\).
    \item Determinar el sentido del campo según el signo de la distribución de carga.
    \item Comprobar las unidades y el comportamiento físico del resultado.
\end{enumerate}


\section{Simetría esférica}

Una distribución posee simetría esférica cuando sus propiedades no cambian al girarla
alrededor de su centro.

En este caso, el campo eléctrico debe ser radial y su magnitud sólo puede depender de
la distancia \(r\) al centro:

\[
E=E(r).
\]

La superficie gaussiana natural es una esfera concéntrica.

Sobre esa esfera:

\[
E=\text{constante}
\]

y el campo es perpendicular a la superficie.

Por tanto,

\[
\Phi_E=E(4\pi r^2).
\]

La ley de Gauss queda

\begin{equation}
\boxed{
E(4\pi r^2)=\frac{Q_{\text{enc}}}{\varepsilon_0}
}.
\label{eq:gauss_simetria_esferica}
\end{equation}


\section{Cáscara esférica uniformemente cargada}

Consideremos una cáscara esférica delgada de radio \(R\) con carga total \(Q\)
distribuida uniformemente sobre su superficie.


\subsection{Campo exterior}

Para

\[
r>R,
\]

una superficie gaussiana esférica de radio \(r\) encierra toda la carga:

\[
Q_{\text{enc}}=Q.
\]

Entonces,

\[
E(4\pi r^2)=\frac{Q}{\varepsilon_0},
\]

y

\begin{equation}
\boxed{
E=\frac{1}{4\pi\varepsilon_0}\frac{|Q|}{r^2}
=k\frac{|Q|}{r^2}
}.
\label{eq:campo_cascara_exterior}
\end{equation}

Exteriormente, una distribución esféricamente simétrica se comporta como si toda su
carga estuviera concentrada en el centro.


\subsection{Campo interior}

Para

\[
r<R,
\]

la superficie gaussiana no encierra carga:

\[
Q_{\text{enc}}=0.
\]

Entonces,

\[
E(4\pi r^2)=0.
\]

Por tanto,

\begin{equation}
\boxed{E=0\qquad(r<R)}.
\label{eq:campo_cascara_interior}
\end{equation}

Este resultado se debe a la simetría esférica y no requiere que la cáscara sea conductora.


\section{Esfera aislante uniformemente cargada}

Consideremos ahora una esfera sólida no conductora de radio \(R\), carga total \(Q\)
y densidad volumétrica uniforme.

La densidad es

\[
\rho=\frac{Q}{\frac43\pi R^3}.
\]


\subsection{Campo exterior}

Para

\[
r\geq R,
\]

la superficie gaussiana encierra toda la carga:

\[
Q_{\text{enc}}=Q.
\]

Por tanto,

\[
\boxed{
E=k\frac{|Q|}{r^2}
\qquad(r\geq R)
}.
\]


\subsection{Campo interior}

Para

\[
r<R,
\]

la superficie gaussiana encierra solamente una parte de la carga.

El volumen encerrado es

\[
V_{\text{enc}}=\frac43\pi r^3.
\]

Por tanto,

\[
Q_{\text{enc}}=\rho V_{\text{enc}}
=\rho\frac43\pi r^3.
\]

Como

\[
\rho=\frac{Q}{\frac43\pi R^3},
\]

obtenemos

\[
Q_{\text{enc}}=Q\frac{r^3}{R^3}.
\]

Aplicamos la ley de Gauss:

\[
E(4\pi r^2)
=
\frac{1}{\varepsilon_0}
Q\frac{r^3}{R^3}.
\]

Despejando,

\[
\boxed{
E=
\frac{|Q|}{4\pi\varepsilon_0R^3}r
}.
\]

También puede escribirse como

\begin{equation}
\boxed{
E=k\frac{|Q|}{R^3}r
\qquad(r<R)
}.
\label{eq:campo_esfera_uniforme_interior}
\end{equation}

Así, dentro de una esfera uniformemente cargada,

\[
\boxed{E\propto r}.
\]

En el centro,

\[
r=0
\Longrightarrow
\boxed{E=0}.
\]

En la superficie,

\[
r=R,
\]

las expresiones interior y exterior dan el mismo valor:

\[
E(R)=k\frac{|Q|}{R^2}.
\]


\subsection{Comportamiento del campo}

Para una esfera sólida uniformemente cargada:

\[
E\propto r
\qquad
(r<R),
\]

mientras que

\[
E\propto\frac{1}{r^2}
\qquad
(r>R).
\]

Por tanto, el campo aumenta linealmente desde cero en el centro hasta alcanzar su
valor máximo en la superficie y luego disminuye con el inverso del cuadrado de la
distancia.


\section{Simetría cilíndrica}

Una línea recta infinitamente larga con carga uniforme posee simetría cilíndrica.

El campo debe ser perpendicular a la línea y su magnitud sólo puede depender de la
distancia perpendicular \(r\):

\[
E=E(r).
\]

Para aprovechar esta simetría elegimos una superficie gaussiana cilíndrica coaxial
con la línea, de radio \(r\) y longitud \(L\).

\begin{center}
\begin{tikzpicture}[>=Latex,scale=0.95]

\draw[dashed] (0,-2.8) -- (0,2.8);
\draw[very thick] (0,-2.5) -- (0,2.5) node[above] {línea cargada};

\draw (-1.6,-2) arc[start angle=180,end angle=360,x radius=1.6,y radius=0.45];
\draw[dashed] (-1.6,-2) arc[start angle=180,end angle=0,x radius=1.6,y radius=0.45];

\draw (-1.6,2) arc[start angle=180,end angle=360,x radius=1.6,y radius=0.45];
\draw[dashed] (-1.6,2) arc[start angle=180,end angle=0,x radius=1.6,y radius=0.45];

\draw (-1.6,-2) -- (-1.6,2);
\draw (1.6,-2) -- (1.6,2);

\draw[->,thick] (1.6,0) -- (2.6,0) node[right] {$E$};
\draw[->,thick] (-1.6,0) -- (-2.6,0) node[left] {$E$};

\node at (0,-3.15) {superficie gaussiana cilíndrica};

\end{tikzpicture}
\end{center}

Sobre la superficie lateral, \(E\) es constante y perpendicular a la superficie.

El área lateral es

\[
A_{\text{lateral}}=2\pi rL.
\]

En las tapas del cilindro el campo es paralelo a las superficies, por lo que

\[
\Phi_{\text{tapas}}=0.
\]

Por tanto,

\[
\Phi_E=E(2\pi rL).
\]


\section{Campo de una línea infinita de carga}

Sea \(\lambda\) la densidad lineal de carga.

La carga encerrada por un cilindro de longitud \(L\) es

\[
Q_{\text{enc}}=\lambda L.
\]

Aplicamos Gauss:

\[
E(2\pi rL)=\frac{\lambda L}{\varepsilon_0}.
\]

Cancelamos \(L\) y despejamos:

\begin{equation}
\boxed{
E=\frac{|\lambda|}{2\pi\varepsilon_0r}
}.
\label{eq:campo_linea_infinita}
\end{equation}

Como

\[
k=\frac{1}{4\pi\varepsilon_0},
\]

también puede escribirse

\[
\boxed{
E=\frac{2k|\lambda|}{r}
}.
\]

Si

\[
\lambda>0,
\]

el campo apunta alejándose radialmente de la línea.

Si

\[
\lambda<0,
\]

apunta hacia ella.

Observe que

\[
\boxed{E\propto\frac{1}{r}}.
\]

El campo de una línea infinita disminuye más lentamente con la distancia que el campo
de una carga puntual.


\section{Simetría plana}

Consideremos una lámina plana, infinita y no conductora con densidad superficial de
carga uniforme \(\sigma\).

Por simetría, el campo debe ser perpendicular al plano.

Además, su magnitud debe ser la misma a ambos lados de la lámina.

Elegimos una superficie gaussiana cilíndrica que atraviese la lámina y cuyas tapas,
de área \(A\), sean paralelas a ella.

Esta superficie suele denominarse \textbf{pastilla gaussiana} o \textbf{pillbox}.

\begin{center}
\begin{tikzpicture}[>=Latex,scale=0.95]

\draw[thick] (-3,0) -- (3,0) node[right] {lámina};

\draw (-1,-1.7) -- (-1,1.7);
\draw (1,-1.7) -- (1,1.7);
\draw (-1,1.7) -- (1,1.7);
\draw (-1,-1.7) -- (1,-1.7);

\draw[->,very thick] (0,1.7) -- (0,2.8) node[above] {$E$};
\draw[->,very thick] (0,-1.7) -- (0,-2.8) node[below] {$E$};

\node at (1.8,1.35) {superficie};
\node at (1.8,0.95) {gaussiana};

\end{tikzpicture}
\end{center}

El campo es perpendicular a las dos tapas y paralelo a la superficie lateral.

Por tanto, la superficie lateral no aporta flujo.

Cada tapa aporta

\[
EA.
\]

El flujo total es

\[
\Phi_E=2EA.
\]


\section{Campo de una lámina infinita uniformemente cargada}

La carga encerrada por la superficie gaussiana es

\[
Q_{\text{enc}}=\sigma A.
\]

La ley de Gauss establece

\[
2EA=\frac{\sigma A}{\varepsilon_0}.
\]

Cancelando \(A\),

\begin{equation}
\boxed{
E=\frac{|\sigma|}{2\varepsilon_0}
}.
\label{eq:campo_lamina_infinita}
\end{equation}

Una característica notable es que la magnitud del campo no depende de la distancia
a la lámina.

Si

\[
\sigma>0,
\]

el campo apunta alejándose de la lámina en ambos lados.

Si

\[
\sigma<0,
\]

apunta hacia ella.


\begin{interpretacion}

La independencia respecto de la distancia es consecuencia de la idealización de una
lámina infinita.

Una lámina física finita sólo se aproxima a este resultado en regiones donde sus
dimensiones sean mucho mayores que la distancia desde el punto de observación hasta
la lámina.

\end{interpretacion}


\section{Comparación de las dependencias con la distancia}

Los resultados anteriores muestran que la geometría de la distribución de carga afecta
la manera en que el campo disminuye con la distancia.

Para una carga puntual o una distribución esféricamente simétrica exterior,

\[
\boxed{E\propto\frac{1}{r^2}}.
\]

Para una línea infinita,

\[
\boxed{E\propto\frac{1}{r}}.
\]

Para una lámina infinita,

\[
\boxed{E=\text{constante respecto de la distancia}}.
\]

Estas diferencias no son casuales. Se relacionan con la manera en que el flujo eléctrico
se distribuye sobre superficies gaussianas cuyo tamaño aumenta.


\section{Cuándo no conviene usar la ley de Gauss para calcular el campo}

La ley de Gauss sigue siendo válida aunque una distribución de carga carezca de
simetría.

Sin embargo, puede ocurrir que no podamos sacar \(E\) fuera de la integral:

\[
\oint_S\vec E\cdot d\vec A.
\]

Por ejemplo, para una distribución irregular de cargas puntuales, el campo puede variar
en magnitud y dirección sobre toda la superficie gaussiana.

En esas situaciones, la ley de Gauss puede proporcionar información sobre el flujo neto,
pero no necesariamente constituye el método más sencillo para calcular el campo en un
punto.


\begin{idea}

\[
\boxed{\text{La ley de Gauss siempre es válida}}
\]

pero

\[
\boxed{\text{no siempre es útil para despejar }E}.
\]

Su mayor poder aparece cuando existe suficiente simetría.

\end{idea}


\section{Ejemplo: flujo a partir de la carga encerrada}

\begin{ejemplo}

Una superficie cerrada encierra tres cargas:

\[
q_1=+\SI{4.0}{\micro\coulomb},\qquad
q_2=-\SI{1.0}{\micro\coulomb},\qquad
q_3=+\SI{2.0}{\micro\coulomb}.
\]

Determine el flujo eléctrico neto a través de la superficie.

\end{ejemplo}

La carga neta encerrada es

\[
Q_{\text{enc}}=q_1+q_2+q_3.
\]

Por tanto,

\[
Q_{\text{enc}}=(4.0-1.0+2.0)\,\si{\micro\coulomb}
=\SI{5.0}{\micro\coulomb}.
\]

Aplicamos la ley de Gauss:

\[
\Phi_E=\frac{Q_{\text{enc}}}{\varepsilon_0}.
\]

Entonces,

\[
\Phi_E=
\frac{5.0\times10^{-6}}{8.854\times10^{-12}}
\approx
\SI{5.6e5}{\newton\metre\squared\per\coulomb}.
\]

Por tanto,

\[
\boxed{
\Phi_E\approx
+\SI{5.6e5}{\newton\metre\squared\per\coulomb}
}.
\]

El signo positivo indica que existe un flujo neto saliente.


\section{Ejemplo: una carga externa}

\begin{ejemplo}

Una superficie cerrada no contiene ninguna carga, pero existe una carga puntual
positiva fuera de ella.

Determine el flujo eléctrico neto a través de la superficie.

\end{ejemplo}

La carga externa produce un campo eléctrico sobre la superficie.

Sin embargo,

\[
Q_{\text{enc}}=0.
\]

Por la ley de Gauss,

\[
\Phi_E=\frac{Q_{\text{enc}}}{\varepsilon_0}=0.
\]

Por tanto,

\[
\boxed{\Phi_E=0}.
\]

Esto no significa que

\[
E=0
\]

en la superficie.

El campo puede ser diferente de cero en todos sus puntos. El flujo neto es cero porque
las contribuciones entrantes y salientes se cancelan.


\section{Densidades de carga y carga encerrada}

En problemas con distribuciones continuas, la ley de Gauss exige calcular
correctamente \(Q_{\text{enc}}\).

Para una distribución lineal,

\[
dq=\lambda\,d\ell,
\]

de modo que

\[
\boxed{
Q_{\text{enc}}=\int\lambda\,d\ell
}.
\]

Para una distribución superficial,

\[
dq=\sigma\,dA,
\]

y

\[
\boxed{
Q_{\text{enc}}=\int\sigma\,dA
}.
\]

Para una distribución volumétrica,

\[
dq=\rho\,dV,
\]

por lo que

\[
\boxed{
Q_{\text{enc}}=\int\rho\,dV
}.
\]

Si la densidad es uniforme, estas expresiones se reducen a

\[
Q_{\text{enc}}=\lambda L,
\]

\[
Q_{\text{enc}}=\sigma A,
\]

o

\[
Q_{\text{enc}}=\rho V,
\]

según el tipo de distribución.


\section{Errores y confusiones frecuentes}

\begin{errorfrecuente}

\textbf{Medir el ángulo con la superficie en lugar de con la normal.}

En

\[
\Phi_E=EA\cos\theta,
\]

\(\theta\) es el ángulo entre el campo y la normal a la superficie.

\end{errorfrecuente}


\begin{errorfrecuente}

\textbf{Pensar que flujo nulo significa campo nulo.}

Puede existir

\[
\Phi_E=0
\]

con

\[
E\neq0.
\]

\end{errorfrecuente}


\begin{errorfrecuente}

\textbf{Incluir cargas externas en \(Q_{\text{enc}}\).}

La ley de Gauss utiliza solamente la carga neta situada dentro de la superficie
gaussiana.

\end{errorfrecuente}


\begin{errorfrecuente}

\textbf{Pensar que las cargas externas no producen campo sobre la superficie.}

Sí pueden producirlo.

Lo que ocurre es que su contribución al flujo neto de una superficie cerrada es cero.

\end{errorfrecuente}


\begin{errorfrecuente}

\textbf{Sumar las magnitudes de las cargas encerradas sin considerar sus signos.}

La carga encerrada es una suma algebraica:

\[
Q_{\text{enc}}=\sum_iq_i.
\]

\end{errorfrecuente}


\begin{errorfrecuente}

\textbf{Utilizar directamente \(\Phi_E=EA\) en cualquier situación.}

Esta expresión sólo se obtiene directamente cuando el campo es constante sobre la
superficie considerada y es perpendicular a ella.

En general debe emplearse

\[
\Phi_E=\int_S\vec E\cdot d\vec A.
\]

\end{errorfrecuente}


\begin{errorfrecuente}

\textbf{Aplicar la ley de Gauss a una superficie abierta.}

La expresión

\[
\oint_S\vec E\cdot d\vec A
=
\frac{Q_{\text{enc}}}{\varepsilon_0}
\]

se aplica a superficies cerradas.

\end{errorfrecuente}


\begin{errorfrecuente}

\textbf{Creer que una superficie gaussiana debe ser una superficie material.}

La superficie gaussiana es imaginaria y se elige para facilitar el análisis.

\end{errorfrecuente}


\begin{errorfrecuente}

\textbf{Suponer que cualquier superficie gaussiana permite calcular fácilmente \(E\).}

La ley de Gauss es siempre válida, pero sólo permite despejar el campo de forma sencilla
cuando la simetría proporciona información suficiente sobre su magnitud y dirección.

\end{errorfrecuente}


\begin{errorfrecuente}

\textbf{Confundir la superficie gaussiana con la distribución de carga.}

La distribución de carga es física.

La superficie gaussiana es una construcción matemática elegida para analizar el campo.

\end{errorfrecuente}


\section{Preguntas para comprobar comprensión}

\begin{enumerate}
    \item ¿Qué representa físicamente el flujo eléctrico?

    \item ¿De qué magnitudes depende el flujo a través de una superficie plana en un
    campo uniforme?

    \item ¿Con respecto a qué dirección se mide el ángulo de
    \[
    \Phi_E=EA\cos\theta?
    \]

    \item ¿Cuándo es máximo el flujo a través de una superficie plana?

    \item ¿Cuándo es cero?

    \item ¿Qué significa que el flujo sea negativo?

    \item ¿Cuál es la unidad SI del flujo eléctrico?

    \item ¿Qué diferencia existe entre una superficie abierta y una cerrada?

    \item ¿En qué dirección se define el vector área de una superficie cerrada?

    \item ¿Qué representa \(Q_{\text{enc}}\)?

    \item ¿Las cargas externas forman parte de \(Q_{\text{enc}}\)?

    \item ¿Puede una carga externa producir campo eléctrico sobre una superficie
    gaussiana?

    \item ¿Por qué una carga externa no produce flujo neto a través de una superficie
    cerrada?

    \item Si \(Q_{\text{enc}}=0\), ¿debe ser necesariamente \(E=0\)?

    \item Si una superficie gaussiana cambia de forma sin dejar salir la carga encerrada,
    ¿cambia el flujo neto?

    \item ¿Qué es una superficie gaussiana?

    \item ¿Por qué la ley de Gauss es especialmente útil en situaciones con simetría?

    \item ¿Qué superficie gaussiana conviene utilizar para una distribución con simetría
    esférica?

    \item ¿Qué superficie conviene para una línea infinita de carga?

    \item ¿Qué superficie conviene para una lámina infinita?

    \item ¿Cómo depende de \(r\) el campo exterior de una distribución esféricamente
    simétrica?

    \item ¿Cómo depende de \(r\) el campo de una línea infinita?

    \item ¿Cómo depende de la distancia el campo de una lámina infinita ideal?

    \item ¿Por qué el campo en el interior de una cáscara esférica uniformemente cargada
    es cero?

    \item ¿Cómo varía el campo dentro de una esfera sólida con carga volumétrica uniforme?
\end{enumerate}


\begin{resumen}

\begin{itemize}
    \item El flujo eléctrico mide la cantidad de campo eléctrico que atraviesa una
    superficie.

    \item Para una superficie plana en un campo uniforme,
    \[
    \boxed{\Phi_E=EA\cos\theta},
    \]
    donde \(\theta\) es el ángulo entre el campo y la normal a la superficie.

    \item Para una superficie general,
    \[
    \boxed{\Phi_E=\int_S\vec E\cdot d\vec A}.
    \]

    \item Para una superficie cerrada,
    \[
    \boxed{\Phi_E=\oint_S\vec E\cdot d\vec A}.
    \]

    \item En una superficie cerrada, el vector área apunta hacia afuera.

    \item El flujo saliente es positivo y el flujo entrante es negativo.

    \item La ley de Gauss es
    \[
    \boxed{
    \oint_S\vec E\cdot d\vec A
    =
    \frac{Q_{\text{enc}}}{\varepsilon_0}
    }.
    \]

    \item La carga encerrada es una suma algebraica:
    \[
    \boxed{Q_{\text{enc}}=\sum_iq_i}.
    \]

    \item Las cargas externas pueden modificar el campo sobre la superficie, pero no
    contribuyen al flujo neto.

    \item Un flujo neto nulo no implica un campo eléctrico nulo.

    \item Una superficie gaussiana es una superficie cerrada imaginaria utilizada para
    aplicar la ley de Gauss.

    \item La ley de Gauss siempre es válida, pero resulta especialmente útil para calcular
    campos con simetría esférica, cilíndrica o plana.

    \item Para una distribución esféricamente simétrica, una superficie gaussiana esférica
    permite escribir
    \[
    \boxed{
    E(4\pi r^2)=\frac{Q_{\text{enc}}}{\varepsilon_0}
    }.
    \]

    \item Fuera de una distribución esféricamente simétrica,
    \[
    \boxed{E=k\frac{|Q|}{r^2}}.
    \]

    \item Dentro de una cáscara esférica uniformemente cargada,
    \[
    \boxed{E=0}.
    \]

    \item Dentro de una esfera sólida uniformemente cargada,
    \[
    \boxed{E=k\frac{|Q|}{R^3}r}.
    \]

    \item Para una línea infinita de carga,
    \[
    \boxed{E=\frac{|\lambda|}{2\pi\varepsilon_0r}}.
    \]

    \item Para una lámina infinita uniformemente cargada,
    \[
    \boxed{E=\frac{|\sigma|}{2\varepsilon_0}}.
    \]

    \item Para distribuciones continuas,
    \[
    Q_{\text{enc}}
    =
    \int\lambda\,d\ell,
    \qquad
    \int\sigma\,dA,
    \qquad
    \int\rho\,dV,
    \]
    según corresponda.
\end{itemize}

\end{resumen}